\documentclass[10pt]{extarticle}
\usepackage[a4paper,margin=1.5cm]{geometry}
\usepackage[T1]{fontenc}
\usepackage{amsmath,amssymb}
\usepackage{graphicx}
\usepackage[export]{adjustbox}
\usepackage{multicol}
\usepackage{enumitem}
\usepackage{booktabs,array,tabularx}
\setlength{\columnsep}{0.7cm}
\setlength{\parindent}{0pt}
\setlength{\parskip}{3pt}
\pagestyle{empty}
\makeatletter
% figure registry (figdefs.tex is generated by the extract step)
\newcommand{\qfigdef}[4]{\@namedef{qf@file@#1}{#2}\@namedef{qf@w@#1}{#3}\@namedef{qf@wide@#1}{#4}}
\newcommand{\qf@check}[1]{\@ifundefined{qf@file@#1}{\PackageError{score2}{Unknown figure id #1}{}}{}}
\newcommand{\qf@img}[4][]{\includegraphics[width=#2\dimexpr\@nameuse{qf@w@#4}pt\relax,max width=#3,max height=0.45\textheight,#1]{\@nameuse{qf@file@#4}}}
% Full-width questions: finalize wraps every question that contains a wide figure in \qfullbegin...\qfullend,
% so the WHOLE question (stem, figures, options) is set across the page in reading order, then the two columns
% resume. The question's head - its number, stem and first figure (\qfullsplit marks the end of the head) -
% is kept on one page, so a number is never left at a page foot with its figure overleaf; the rest may break
% across pages normally, so no big gaps are left. Without a \qfullsplit, the whole question is kept together
% if it is under 45% of a page. Never write these yourself.
\newif\ifqf@full
\newif\ifqf@boxed
\newsavebox{\qf@box}
\newcommand{\qf@open}{\global\qf@fulltrue\par\global\qf@boxedtrue
  \setbox\qf@box\vbox\bgroup\hsize\textwidth\linewidth\textwidth}
\newcommand{\qf@emit}[1]{\par\egroup\global\qf@boxedfalse
  \ifdim\dimexpr\ht\qf@box+\dp\qf@box\relax<#1\textheight\noindent\usebox{\qf@box}\par\else\unvbox\qf@box\fi}
\newcommand{\qfullsplit}{\ifqf@boxed\qf@emit{0.85}\fi}
\newcommand{\qf@close}{\ifqf@boxed\qf@emit{0.45}\fi\global\qf@fullfalse}
\newcommand{\qfullbegin}{\end{multicols}\qf@open}
\newcommand{\qfullend}{\qf@close\begin{multicols}{2}}
\newcommand{\qfullstart}{\qf@open}
\newcommand{\qfulldone}{\qf@close}
% \qfig[scale]{id}: block figure, centred, directly where it is written (size set by the tool).
% Fallback if a wide figure is not inside a full-width question: break out of the columns in place - no floats.
\newcommand{\qfig}[2][1]{\qf@check{#2}%
  \ifqf@full\qf@blk{#1}{#2}\else
  \ifnum\@nameuse{qf@wide@#2}=1 \expandafter\@firstoftwo\else\expandafter\@secondoftwo\fi
  {\end{multicols}\begin{center}\qf@img{#1}{\textwidth}{#2}\end{center}\begin{multicols}{2}}%
  {\qf@blk{#1}{#2}}\fi}
\newcommand{\qf@blk}[2]{\par\smallskip{\centering\qf@img{#1}{\linewidth}{#2}\par}\smallskip}
% \qfigi[scale]{id}: inline figure (inside option lists / tables / minipages); never breaks columns.
\newcommand{\qfigi}[2][1]{\qf@check{#2}\qf@img[valign=c]{#1}{\linewidth}{#2}}
\makeatother
% structure macros - lint reads these to check completeness and figure placement
\newcommand{\papersec}[1]{\par\bigskip{\centering\Large\bfseries PAPER-#1\par}\medskip}
\newcommand{\partsec}[2]{\par\medskip\noindent{\large\bfseries PART-#1: #2}\par\smallskip}
\newcommand{\sect}[1]{\par\smallskip\noindent\textbf{#1}\par}
\newcommand{\qn}[1]{\par\smallskip\noindent\textbf{#1.}~}
\newcommand{\paraq}[2]{\par\smallskip\noindent\textbf{Paragraph for Questions #1 and #2}\par}
\newenvironment{qoptions}{\begin{itemize}[leftmargin=1.8em,itemsep=1pt,topsep=2pt,label={}]}{\end{itemize}}
\newcommand{\opt}[2]{\item[(#1)] #2}
% figure registry (generated by score2_tools.py finalize)
\qfigdef{p1_phy_s1_q01_1}{TMAS - Test 04 (19 Nov 2023)-img/p1_phy_s1_q01_1.png}{159.9}{0}
\qfigdef{p1_phy_s1_q03_1}{TMAS - Test 04 (19 Nov 2023)-img/p1_phy_s1_q03_1.png}{77.9}{0}
\qfigdef{p1_phy_s1_q03_A}{TMAS - Test 04 (19 Nov 2023)-img/p1_phy_s1_q03_A.png}{70.8}{0}
\qfigdef{p1_phy_s1_q03_B}{TMAS - Test 04 (19 Nov 2023)-img/p1_phy_s1_q03_B.png}{71.0}{0}
\qfigdef{p1_phy_s1_q03_C}{TMAS - Test 04 (19 Nov 2023)-img/p1_phy_s1_q03_C.png}{75.6}{0}
\qfigdef{p1_phy_s1_q03_D}{TMAS - Test 04 (19 Nov 2023)-img/p1_phy_s1_q03_D.png}{71.0}{0}
\qfigdef{p1_phy_s2_q12_1}{TMAS - Test 04 (19 Nov 2023)-img/p1_phy_s2_q12_1.png}{239.2}{0}
\qfigdef{p1_phy_s3_q02_1}{TMAS - Test 04 (19 Nov 2023)-img/p1_phy_s3_q02_1.png}{169.7}{0}
\qfigdef{p1_phy_s3_q06_1}{TMAS - Test 04 (19 Nov 2023)-img/p1_phy_s3_q06_1.png}{222.4}{0}
\qfigdef{p1_chem_s2_q08_1}{TMAS - Test 04 (19 Nov 2023)-img/p1_chem_s2_q08_1.png}{26.5}{0}
\qfigdef{p1_chem_s2_q08_2}{TMAS - Test 04 (19 Nov 2023)-img/p1_chem_s2_q08_2.png}{58.0}{0}
\qfigdef{p1_chem_s2_q08_B}{TMAS - Test 04 (19 Nov 2023)-img/p1_chem_s2_q08_B.png}{44.8}{0}
\qfigdef{p1_chem_s2_q08_4}{TMAS - Test 04 (19 Nov 2023)-img/p1_chem_s2_q08_4.png}{43.8}{0}
\qfigdef{p1_chem_s2_q08_C}{TMAS - Test 04 (19 Nov 2023)-img/p1_chem_s2_q08_C.png}{49.7}{0}
\qfigdef{p1_chem_s2_q08_6}{TMAS - Test 04 (19 Nov 2023)-img/p1_chem_s2_q08_6.png}{51.5}{0}
\qfigdef{p1_chem_s2_q08_7}{TMAS - Test 04 (19 Nov 2023)-img/p1_chem_s2_q08_7.png}{100.9}{0}
\qfigdef{p1_chem_s3_q02_1}{TMAS - Test 04 (19 Nov 2023)-img/p1_chem_s3_q02_1.png}{109.0}{0}
\qfigdef{p1_chem_s3_q04_1}{TMAS - Test 04 (19 Nov 2023)-img/p1_chem_s3_q04_1.png}{154.3}{0}
\qfigdef{p1_chem_s3_q05_1}{TMAS - Test 04 (19 Nov 2023)-img/p1_chem_s3_q05_1.png}{147.7}{0}
\qfigdef{p2_phy_s1_q04_1}{TMAS - Test 04 (19 Nov 2023)-img/p2_phy_s1_q04_1.png}{100.1}{0}
\qfigdef{p2_phy_s2_q01_1}{TMAS - Test 04 (19 Nov 2023)-img/p2_phy_s2_q01_1.png}{150.1}{0}
\qfigdef{p2_phy_s2_q03_1}{TMAS - Test 04 (19 Nov 2023)-img/p2_phy_s2_q03_1.png}{160.7}{0}
\qfigdef{p2_phy_s2_q04_1}{TMAS - Test 04 (19 Nov 2023)-img/p2_phy_s2_q04_1.png}{219.4}{0}
\qfigdef{p2_phy_s2_q05_1}{TMAS - Test 04 (19 Nov 2023)-img/p2_phy_s2_q05_1.png}{58.3}{0}
\qfigdef{p2_phy_s3_q05_1}{TMAS - Test 04 (19 Nov 2023)-img/p2_phy_s3_q05_1.png}{332.1}{1}
\qfigdef{p2_chem_s1_q03_1}{TMAS - Test 04 (19 Nov 2023)-img/p2_chem_s1_q03_1.png}{182.5}{0}
\raggedcolumns
\makeatletter\@beginparpenalty=10000\makeatother
\newcolumntype{L}[1]{>{\raggedright\arraybackslash}p{#1}}
\begin{document}
\begin{center}{\Large\bfseries TMAS --- Test 04 (19 Nov 2023)}\end{center}
\begin{multicols}{2}
\raggedcolumns
\papersec{1}
\partsec{1}{PHYSICS}
\sect{SECTION-I(i)}
\qn{1} In a convex lens of focal length $10$ cm, a small portion of length $2$ mm is removed symmetrically and the two parts are joined together. A source is placed symmetrically $5$ cm from composite system and a screen is placed $120$ cm from composite system.
\qfig{p1_phy_s1_q01_1}
\begin{qoptions}
\opt{A}{No interference fringes can be seen on screen}
\opt{B}{To see interference fringes, screen should be moved towards composite system.}
\opt{C}{To see interference fringes, source should be moved away from composite system.}
\opt{D}{Interference fringes will be obtained in region of width $24$ cm.}
\end{qoptions}
\qn{2} In an experiment to study standing waves, you use a string whose mass per length is $\mu=(1.0\pm0.1)\times10^{-4}$ kg/m. You look at the fundamental mode, whose frequency $f$ is related to the length $L$ and tension $T$ of the string by the following equation $L=\dfrac{1}{2f}\sqrt{\dfrac{T}{\mu}}$. You make a plot with $L$ on the $y$-axis and $\sqrt T$ on the $x$-axis, and find that the best fitting line is $y=(8.0\pm0.3)\times10^{-3}x+(0.2\pm0.04)$ in SI units. What is the value of the frequency of the wave (including the error)? Express your result in SI unit (Hz).
\begin{qoptions}
\opt{A}{$6250\pm859.8$ Hz}
\opt{B}{$1250\pm32.3$ Hz}
\opt{C}{$6250\pm546.9$ Hz}
\opt{D}{$4875\pm287$ Hz}
\end{qoptions}
\qn{3} An open pipe of sufficient length starts to dip in water at $t=0$ with velocity $v$. If sound velocity in air is $v_0$ then correct graph between rate of change of fundamental frequency $\dfrac{df_0}{dt}$ and $t$ is
\qfig{p1_phy_s1_q03_1}
\begin{qoptions}
\opt{A}{\qfigi{p1_phy_s1_q03_A}}
\opt{B}{\qfigi{p1_phy_s1_q03_B}}
\opt{C}{\qfigi{p1_phy_s1_q03_C}}
\opt{D}{\qfigi{p1_phy_s1_q03_D}}
\end{qoptions}
\qn{4} A bird is singing on a tree and a man is hearing at a distance `$r$' from the bird. Calculate the displacement of the man towards the bird so that the loudness heard by man increases by $20$ dB.

[Assume that the motion of man is along the line joining the bird and the man]
\begin{qoptions}
\opt{A}{$\dfrac{9r}{10}$}
\opt{B}{$\dfrac{r}{10}$}
\opt{C}{$\dfrac{3r}{5}$}
\opt{D}{$\dfrac{4r}{5}$}
\end{qoptions}
\qn{5} The radius of the orbit of an electron in a Hydrogen-like atom is $4.5\,a_0$, where $a_0$ is the Bohr radius. Its orbital angular momentum is $\dfrac{3h}{2\pi}$. It is given that $h$ is Planck constant and $R$ is Rydberg constant. The wavelength, this atom would not be emitting, as it de-excites itself is
\begin{qoptions}
\opt{A}{$\dfrac{9}{32R}$}
\opt{B}{$\dfrac{9}{16R}$}
\opt{C}{$\dfrac{9}{5R}$}
\opt{D}{$\dfrac{1}{3R}$}
\end{qoptions}
\qn{6} The radionuclide ${}^{11}_{6}\mathrm{C}$ decays by $\beta^+$ emission.

Given that $m({}^{11}_{6}\mathrm{C})=11.011434$ u

$m({}^{11}_{5}\mathrm{B})=11.009305$ u

$m_e=0.000548$ u, $1\,\mathrm{u}=931.5$ MeV/$c^2$

The Q-value of this decay process is :-
\begin{qoptions}
\opt{A}{$0.962$ MeV}
\opt{B}{$0.962\times10^{3}$ MeV}
\opt{C}{$0.962$ eV}
\opt{D}{$0$}
\end{qoptions}
\sect{SECTION-I(ii)}
\qn{7} Mark the correct options-
\begin{qoptions}
\opt{A}{An atom with a vacancy has smaller energy than a neutral atom}
\opt{B}{$K_\alpha$ X-ray is emitted when a hole makes a jump from K-shell to some other shell}
\opt{C}{The wavelength of K X-ray is smaller than the wavelength of L X-ray of the same material}
\opt{D}{The wavelength of $K_\alpha$ X-ray is smaller than the wavelength of $K_\beta$ X-ray of the same material}
\end{qoptions}
\qn{8} In an X-ray tube, electrons emitted from a filament (cathode) carrying current $I$ hit a target (anode) at a distance $d$ from the cathode. The target is kept at a potential $V$ higher than the cathode resulting in emission of continuous and characteristic X-rays. If the filament current $I$ is decreased to $\dfrac I2$, the potential difference $V$ is increased to $2V$ and the separation distance $d$ is reduced to $\dfrac d2$, then :-
\begin{qoptions}
\opt{A}{the cut-off wavelength will reduce to half and the wavelengths of the characteristic X-rays will remain the same.}
\opt{B}{the cut-off wavelength as well as the wavelengths of the characteristic X-rays will remain the same.}
\opt{C}{the cut-off wavelength will reduce to half and the intensities of all the X-rays will decrease.}
\opt{D}{the cut-off wavelength will become two times larger, and the intensity of all the X-rays will decrease.}
\end{qoptions}
\qn{9} In $\beta$-decay, the $Q$-value of the process is $E$. Then
\begin{qoptions}
\opt{A}{$K.E.$ of a $\beta$-particle cannot exceed $E$.}
\opt{B}{$K.E.$ of anti neutrino emitted lies between Zero and $E$.}
\opt{C}{$N/Z$ ratio of the nucleus is altered.}
\opt{D}{Mass number ($A$) of the nucleus is altered.}
\end{qoptions}
\qn{10} The pitch of screw gauge is $1$ mm and its circular scale is divided into $100$ divisions. When nothing is put between the studs the zero of main scale is not seen, but when circular scale is rotated by $450^\circ$ the zero of main scale is seen and the zero of main scale coincide with the zero of circular scale. When a glass plate is placed between the studs, the circular scale lies between $18^{\mathrm{th}}$ to $19^{\mathrm{th}}$ division of main scale and circular scale reads $34$ divisions. Then,
\begin{qoptions}
\opt{A}{Error is positive zero error and its magnitude is $1.25$ mm.}
\opt{B}{Error is negative zero error and its magnitude is $1.25$ mm.}
\opt{C}{The thickness of the glass plate is $19.59$ mm.}
\opt{D}{The thickness of the glass plate is $17.09$ mm.}
\end{qoptions}
\qn{11} The period of oscillation of a simple pendulum in an experiment is recorded as $2.63$ sec, $2.56$ sec, $2.42$ sec, $2.71$ sec and $2.80$ sec respectively:-
\begin{qoptions}
\opt{A}{Time period is $2.62$ sec}
\opt{B}{Average absolute error is $0.01$ sec}
\opt{C}{Percentage error is $2.4\%$}
\opt{D}{Percentage error is $4.2\%$}
\end{qoptions}
\qn{12} In given young's double slit experiment a point sources of wavelength $\lambda=500$ nm is placed slightly off the central axis. Point A and B are in front of upper and lower slit respectively. If intensity from upper \& lower slits is $I_0$ each on the screen, then
\qfig{p1_phy_s2_q12_1}
\begin{qoptions}
\opt{A}{Intensity at A, O and B is $4I_0$, $2I_0$ and $4I_0$ respectively.}
\opt{B}{$70^{\mathrm{th}}$ maxima is formed at B}
\opt{C}{$30^{\mathrm{th}}$ maxima is formed at A}
\opt{D}{If a glass slab of thickness $0.02$ mm and refractive index $\mu=3/2$ is kept in front of lower slit then central maxima is formed exactly at O.}
\end{qoptions}
\sect{SECTION-II}
\qn{1} In a single slit diffraction experiment first minimum of red light ($660$ nm) coincides with first maximum of some other wavelength (in nm) $\lambda$. Calculate $\lambda$.
\qn{2} A glass lens is coated on one side with a thin film of magnesium fluoride ($\mathrm{MgF_2}$) to reduce reflection from the lens surface. The index of refraction of $\mathrm{MgF_2}$ is $1.38$; that of the glass is $1.50$. What is the least coating thickness that eliminates (via interference) the reflections at the middle of the visible spectrum ($\lambda=550$ nm) ? Assume that the light is approximately perpendicular to the lens surface. Express your answer in nm.
\qfig{p1_phy_s3_q02_1}
\qn{3} The period of oscillation of a simple pendulum is $T=2\pi\sqrt{L/g}$. The measured value of $L$ is $20.0$ cm known to $1$ mm accuracy and time for $100$ oscillations of the pendulum is found to be $90$ s using a wrist-watch of $1$ s resolution. If the accuracy in the determination of $g$ is $n\%$, then the value of $2n$ is (in rounded off to nearest integer)
\qn{4} The diameter of a cylinder is measured using a Vernier callipers with no zero error. It is found that the zero of the Vernier scale lies between $5.10$ cm and $5.15$ cm of the main scale. The Vernier scale has $50$ divisions which coincide with $49$ main scale divisions. The $24^{\mathrm{th}}$ division of the Vernier scale exactly coincides with one of the main scale divisions. The diameter of the cylinder is (in mm) :
\qn{5} The circumference of circular orbit of electron in a $\mathrm{He^+}$ ion is five times the de-Broglie wavelength associated with the electron. Find the radius of the orbit.

(Express answer in \AA\ and take Bohr radius of hydrogen as $0.529$ \AA)
\qn{6} A parallel beam of electrons travelling in $x$--direction falls on a slit AB of width $d$ as shown in the figure. After passing through the slit, an electron acquires momentum $p_y$ in the $y$--direction and hit the screen at distance $\dfrac{3\lambda D}{4d}$ from $O$. If $h$ be the planks constant then find the value of $\dfrac{5p_yd}{3h}$.
\qfig{p1_phy_s3_q06_1}
Here $O$ is a point equidistant from A \& B
\partsec{2}{CHEMISTRY}
\sect{SECTION-I(i)}
\qn{1} Which one of the following statements regarding helium is INCORRECT?
\begin{qoptions}
\opt{A}{It is used to produce and sustain powerful superconducting magnets.}
\opt{B}{It is used as a cryogenic agent for carrying out experiments at low temperatures.}
\opt{C}{It is used in filling balloons for meterological observations because it is lighter and non-inflammable.}
\opt{D}{It is used as dilute for oxygen in modern diving apparatus because of its high solubility in blood.}
\end{qoptions}
\qn{2} Statement--1 : $\mathrm{Mn_2O_7}$ is a covalent green oil.

Statement--2 : CrO is basic but $\mathrm{Cr_2O_3}$ is amphoteric

Statement--3 : $\mathrm{FeO_4^{2-}}$ formed in alkaline medium

Assign True (T) and False (F) for above statements.
\begin{qoptions}
\opt{A}{FTT}
\opt{B}{FFT}
\opt{C}{TTT}
\opt{D}{FFF}
\end{qoptions}
\qn{3} The enthalpy of neutralisation of four acids HA, HB, HC and HD with KOH are $-10.2$, $-11.4$, $-12.3$ and $-13.1$ kcal/mole respectively. Which of the following aqueous salt solutions has maximum pH ?
\begin{qoptions}
\opt{A}{$0.1$ M NaA}
\opt{B}{$0.1$ M NaB}
\opt{C}{$0.1$ M NaC}
\opt{D}{$0.1$ M NaD}
\end{qoptions}
\qn{4} Which of the following option(s) is/are correct regarding point defects in solids?
\begin{qoptions}
\opt{A}{All point defects result in decreasing the density of solid.}
\opt{B}{The density of point defect per unit volume of solid increases with decrease in temperature.}
\opt{C}{In the formation of impurity defects, density of solid increases always.}
\opt{D}{In the formation of point defects, the entropy of system increases but the entropy of surrounding decreases.}
\end{qoptions}
\qn{5} Select the correct statement(s) from the following option :
\begin{qoptions}
\opt{A}{If water boils at temperature $T^\circ$C at an external pressure of $2$ atm then $[T>100^\circ\mathrm{C}]$}
\opt{B}{For a hypothetical gaseous reaction $\mathrm{A(g)}\rightleftharpoons3\mathrm{B(g)}$, at $10$ Kelvin, $K_P>K_C$ (P in terms of atm and C in terms of mol/litre)}
\opt{C}{Freezing point of water at an external pressure of $2$ atm is greater than its normal freezing point.}
\opt{D}{Position of equilibrium in a reaction, $\mathrm{NH_2COONH_4(s)}\rightleftharpoons2\mathrm{NH_3(g)}+\mathrm{CO_2(g)}$ at a given temperature remains unchanged on addition of more $\mathrm{NH_3(g)}$ at equilibrium.}
\end{qoptions}
\qn{6} Which of the following statement is INCORRECT?
\begin{qoptions}
\opt{A}{Terfenadine has `$3$' aromatic rings in it's structure}
\opt{B}{Prontosil is coloured due to conjugated double bond system.}
\opt{C}{Chloromphenicol contains phenolic functional group.}
\opt{D}{Chloroxylenol has a phenolic functional group}
\end{qoptions}
\sect{SECTION-I(ii)}
\qn{7} Which of the following statements is / are CORRECT for $\mathrm{I_2O_5}$ :
\begin{qoptions}
\opt{A}{It is oxidising agent}
\opt{B}{It is used to estimate CO in a gas mixture}
\opt{C}{It is anhydride of $\mathrm{HIO_4}$ (periodic acid)}
\opt{D}{When it react with mixture of $\mathrm{CO+CO_2}$ all the $\mathrm{CO_2}$ is converted to CO.}
\end{qoptions}
\qn{8} Reaction incorrectly matched with major product :
\begin{qoptions}
\opt{A}{\qfigi{p1_chem_s2_q08_1}\ $\xrightarrow[\text{2. }\mathrm{H_2O}/\Delta]{\text{1. }\mathrm{O_3}}$\ \qfigi{p1_chem_s2_q08_2}}
\opt{B}{\qfigi{p1_chem_s2_q08_B}\ $\xrightarrow{\mathrm{Br_2/H_2O}}$\ \qfigi{p1_chem_s2_q08_4}}
\opt{C}{\qfigi{p1_chem_s2_q08_C}\ \qfigi{p1_chem_s2_q08_6}\par\qfigi{p1_chem_s2_q08_7}}
\opt{D}{{\small$\mathrm{CH_3CH_2CH_2COOH}\xrightarrow{\text{Red P}+\mathrm{Cl_2}}\mathrm{CH_3CH_2CH(Cl)COOH}$}}
\end{qoptions}
\qn{9} 
\begin{center}
\begin{tabular}{|l|l|l|l|}\hline
\multicolumn{2}{|c|}{\textbf{List-I}} & \multicolumn{2}{c|}{\textbf{List-II}}\\\hline
(a) & Vitamin $\mathrm{B_2}$ & (P) & Ascorbic acid\\\hline
(b) & Vitamin C & (Q) & Artificial sweetner\\\hline
(c) & Sucralose & (R) & Thiamine\\\hline
(d) & Vitamin $\mathrm{B_1}$ & (S) & Riboflavin\\\hline
\end{tabular}
\end{center}
The incorrect match is/are?
\begin{qoptions}
\opt{A}{a$\to$P ; b$\to$Q ; c$\to$S ; d$\to$R}
\opt{B}{a$\to$S ; b$\to$P ; c$\to$Q ; d$\to$R}
\opt{C}{a$\to$Q ; b$\to$P ; c$\to$R ; d$\to$S}
\opt{D}{a$\to$R ; b$\to$Q ; c$\to$P ; d$\to$S}
\end{qoptions}
\qn{10} Select the incorrect statement(s) about compounds (X) and (Y) formed in below given reaction.

$\mathrm{P_4+SO_2Cl_2}\to\underset{\text{(planar molecule)}}{\mathrm{X}}+\underset{\text{(Non planar molecule)}}{\mathrm{Y}}$
\begin{qoptions}
\opt{A}{Both X and Y are polar molecules in gaseous state}
\opt{B}{Both X and Y are non polar molecules in gaseous state}
\opt{C}{In both X and Y, $3d$ orbitals used in bonding}
\opt{D}{Complete hydrolysis of both X and Y gives dibasic oxy acid}
\end{qoptions}
\qn{11} Which of the following compound(s) has peroxy linkage ?
\begin{qoptions}
\opt{A}{Caro's acid}
\opt{B}{Marshall's acid}
\opt{C}{Perchloric acid}
\opt{D}{Iodic acid}
\end{qoptions}
\qn{12} Among the following, the diamagnetic compound(s) is / are
\begin{qoptions}
\opt{A}{$\mathrm{Na_2O_2}$}
\opt{B}{$\mathrm{O_3}$}
\opt{C}{$\mathrm{KO_2}$}
\opt{D}{$\mathrm{N_2O}$}
\end{qoptions}
\sect{SECTION-II}
\qn{1} Two moles of an ideal solution is prepared by mixing pure liquids `A' and `B' at $27^\circ$C. The mole percent of `A' in the vapour above the solution is $60\%$. Vapour pressure of pure liquids A and B are $120$ mmHg and $80$ mmHg respectively. The enthalpy change of mixing $|\Delta_{\mathrm{mix}}H|$ of the solution is \ldots\ldots\ldots calories.

(Given : $R=2$ cal K/mol, $\log2=0.3$, $\log3=0.48$, $\log5=0.7$, $\ln10=2.3$)
\qn{2} Rate constant ($k$) of elementary reaction $\mathrm{A(g)}\to\mathrm{B(g)}$ at $27^\circ$C is \underline{\hspace{1cm}}$\times10^{7}$ (sec$^{-1}$). Following plot is given for this reaction with units of $k$ as sec$^{-1}$.
\qfig{p1_chem_s3_q02_1}
\qn{3} Given :

$[\mathrm{Ag(NH_3)_2}^+](\mathrm{aq})\rightleftharpoons\mathrm{Ag}^+(\mathrm{aq})+2\mathrm{NH_3(aq)}$, $K_C=6\times10^{-7}$ and $K_{sp}$ of AgCl $=2.4\times10^{-10}$ at $298$ K

The concentration of the complex $[\mathrm{Ag(NH_3)_2}^+]$ in resultant solution, when solid AgCl is added to aqueous solution having $1$M ammonia to form saturated solution is $P\times10^{-3}$ M. Find value of $P$ ?
\qn{4} The structure of a peptide is given below.
\qfig{p1_chem_s3_q04_1}
If the absolute values of the net charge of the peptide at pH $=2$, pH $=6$, and pH $=11$ are $|z_1|$, $|z_2|$ and $|z_3|$, respectively, then what is $|z_1|+|z_2|+|z_3|$ ?
\qn{5} 
\qfig{p1_chem_s3_q05_1}
At which position nitration mainly takes place ?
\qn{6} How many of the following drugs is/are example(s) of broad spectrum antibiotic ?

Ofloxacin, Penicillin--G, Terpineol, Salvarsan
\partsec{3}{MATHEMATICS}
\sect{SECTION-I(i)}
\qn{1} Number of solutions of the equation $\left|\tan^{-1}|x|\right|=\sqrt{\left(x^2+\dfrac12\right)^2-2x^2}$ is
\begin{qoptions}
\opt{A}{$1$}
\opt{B}{$2$}
\opt{C}{$3$}
\opt{D}{$4$}
\end{qoptions}
\qn{2} Let $f(x)$ be a twice differentiable function satisfying $xf(x^2)+(x^2+1)f(x)=x,\ \forall\ x\in R$ then the value of $f''(1)$ is
\begin{qoptions}
\opt{A}{$\dfrac13$}
\opt{B}{$-\dfrac13$}
\opt{C}{$-\dfrac19$}
\opt{D}{$1$}
\end{qoptions}
\qn{3} Suppose that $f$ is differentiable function such that $f'(x)\le2$ for all $x$. If $f(1)=2$ and $f(4)=8$ then $f(2)$ has the value equal to
\begin{qoptions}
\opt{A}{$4$}
\opt{B}{$2$}
\opt{C}{$1$}
\opt{D}{$0$}
\end{qoptions}
\qn{4} Let $\displaystyle\int\dfrac{x^2+20}{(x\sin x+5\cos x)^2}\,dx=f(x)+\tan x+C$, $f(0)=0$ where $C$ denotes constant of integration. Then the value of $\pi+5f(\pi)$ is equal to
\begin{qoptions}
\opt{A}{$2$}
\opt{B}{$4$}
\opt{C}{$3$}
\opt{D}{$0$}
\end{qoptions}
\qn{5} Number of values of $x$ satisfying the system of equation $\sin^{-1}\!\left(\dfrac{2x}{1+x^2}\right)=\cos^{-1}\!\left(\dfrac{1-x^2}{1+x^2}\right)=2\cos^{-1}x$ is
\begin{qoptions}
\opt{A}{$0$}
\opt{B}{$1$}
\opt{C}{$2$}
\opt{D}{infinite}
\end{qoptions}
\qn{6} The solution of $\dfrac{dy}{dx}=\dfrac{x^2+y^2+1}{2xy}$ satisfying $y(1)=1$ is given by :
\begin{qoptions}
\opt{A}{a system of hyperbola}
\opt{B}{a system of circles}
\opt{C}{$y^2=x(1+x)-1$}
\opt{D}{$(x-2)^2+(y-3)^2=5$}
\end{qoptions}
\sect{SECTION-I(ii)}
\qn{7} The value of $\displaystyle\sum_{r=2}^{100}\dfrac{2^r(2-r)}{r(r+1)(r+2)}$ is equal to
\begin{qoptions}
\opt{A}{$\dfrac23-\dfrac{2^{100}}{(51)(101)}$}
\opt{B}{$\dfrac23-\dfrac{2^{101}}{(101)(102)}$}
\opt{C}{$\dfrac12-\dfrac{2^{100}}{100.101}$}
\opt{D}{$\dfrac12-\dfrac{2^{100}}{(101)(51)}$}
\end{qoptions}
\qn{8} Let $\vec A=\hat i-2\hat j+3\hat k$, $\vec B=2\hat i+\hat j-\hat k$, $\vec C=\hat j+\hat k$. If the vector $\vec B\times\vec C$ can be expressed as $\vec B\times\vec C=x\vec A+y\vec B+z\vec C$ where $x,y,z$ are scalars, then which of the following is/are correct
\begin{qoptions}
\opt{A}{$x+2z=0$}
\opt{B}{$x=2y=2z=1$}
\opt{C}{$y+z=0$}
\opt{D}{$x=2y$}
\end{qoptions}
\qn{9} Let P denotes the plane consisting of all points that are equidistant from the points $A(-4,2,1)$ and $B(2,-4,3)$ and Q be the plane $x-y+cz=1$ where $c\in R$.
\begin{qoptions}
\opt{A}{The plane P is parallel to plane Q if $c=1/3$}
\opt{B}{The plane P is parallel to plane Q if $c=3$}
\opt{C}{If angle between the planes P \& Q is $45^\circ$, then product of all possible values of `$c$' is $-2$}
\opt{D}{If angle between the planes P \& Q is $45^\circ$, then sum of all possible values of `$c$' is $\dfrac{24}{17}$}
\end{qoptions}
\qn{10} If A and B are two matrices such that $AB=A$ and $BA=B$, then
\begin{qoptions}
\opt{A}{$(A^3-B^3)^5=(A^5-B^5)^3$}
\opt{B}{$(A^4-B^4)^3=A^{12}-B^{12}$}
\opt{C}{$(A-B)$ is idempotent}
\opt{D}{$(A^{10}-B^{10})$ is nilpotent}
\end{qoptions}
\qn{11} Which of the following option is/are correct ?
\begin{qoptions}
\opt{A}{If $\vec a=\alpha\hat i+2\hat j-3\hat k$, $\vec b=\hat i+2\alpha\hat j-2\hat k$ and $\vec c=2\hat i-\alpha\hat j+\hat k$ then $\left\{(\vec a\times\vec b)\times(\vec b\times\vec c)\right\}\times(\vec c\times\vec a)=0$ if $\alpha=\dfrac23$}
\opt{B}{If $\vec a=\alpha\hat i+2\hat j-3\hat k$, $\vec b=\hat i+2\alpha\hat j-2\hat k$ and $\vec c=2\hat i-\alpha\hat j+\hat k$ then $\left\{(\vec a\times\vec b)\times(\vec b\times\vec c)\right\}\times(\vec c\times\vec a)=0$ if $\alpha=\dfrac32$}
\opt{C}{If $\lambda\hat i+\hat j+2\hat k$, $\hat i+\lambda\hat j-\hat k$ and $2\hat i-\hat j+\lambda\hat k$ are coplanar then $\lambda$ is $-2$}
\opt{D}{If $\lambda\hat i+\hat j+2\hat k$, $\hat i+\lambda\hat j-\hat k$ and $2\hat i-\hat j+\lambda\hat k$ are coplanar then sum of all possible values of $\lambda$ is $0$}
\end{qoptions}
\qn{12} Let $A_n=\displaystyle\sum_{k=0}^{n-1}\dfrac{n}{k^2+n^2+(n+2)k+n+1}$ and $B_n=\displaystyle\sum_{k=0}^{n-1}\dfrac{n}{n^2+kn+k^2}$ for $n=1,2,3,\ldots\ldots$, then which of the following is/are true ?
\begin{qoptions}
\opt{A}{$A_n<B_n$}
\opt{B}{$A_n>B_n$}
\opt{C}{$A_n<\dfrac\pi3$}
\opt{D}{$B_n>\dfrac\pi9$}
\end{qoptions}
\sect{SECTION-II}
\qn{1} Let $A=\{1,2,3,4\}$, $B=\{1,2,3,4\}$. The total number of functions $f:A\to B$ such that $f(x)\neq x\ \forall\ x\in A$ and range of $f(x)$ contains exactly three elements, is
\qn{2} Consider two lines $L_1:\dfrac{x-7}{3}=\dfrac{y-7}{2}=\dfrac{z-3}{1}$ and $L_2:\dfrac{x-1}{2}=\dfrac{y+1}{4}=\dfrac{z+1}{3}$. If a line $L$ whose direction ratios are $<2,2,1>$ intersect the line $L_1$ and $L_2$ at A and B, then find the distance AB.
\qn{3} Let $A_1,A_2,A_3,A_4$ be the areas of the triangular faces of a tetrahedron and $h_1,h_2,h_3,h_4$ be the corresponding altitudes drawn from the opposite vertices of the tetrahedron. If volume of tetrahedron is $5$ cubic units, then find the minimum value of $(A_1+A_2+A_3+A_4)(h_1+h_2+h_3+h_4)$ (in cubic units).
\qn{4} Let ABCD be any quadrilateral. P and Q are the midpoints of its diagonal.

If $\overrightarrow{AB}^2+\overrightarrow{BC}^2+\overrightarrow{CD}^2+\overrightarrow{DA}^2-\overrightarrow{AC}^2-\overrightarrow{BD}^2=\lambda\overrightarrow{PQ}^2$, then find the value of $\lambda$.
\qn{5} The value of $\displaystyle\int_0^1\sqrt[3]{2x^3-3x^2-x+1}\,dx$ is -
\qn{6} Given $f$ is a differentiable function such that $1\leqslant f(x)+f'(x)\leqslant2\ \forall\ x\in R$ and $f(0)=0$, then the value of $\left[\dfrac{f(1)}{2}\right]$ is equal to (where $[x]$ denotes the greatest integer less than or equal to real number $x$)
\raggedcolumns
\papersec{2}
\partsec{1}{PHYSICS}
\sect{SECTION-I}
\qn{1} $Q$ is energy released in the decy, $m_x$ is atomic mass of pareant nucleus, $m_y$ is atomic mass of daughter nucleus and $m_e$ is mass of electron, then match the following :-
\begin{qoptions}
\opt{A}{For k-capture $Q=(m_x-m_y)c^2$}
\opt{B}{For $\beta^-$ decay $Q=(m_x-m_y)c^2$}
\opt{C}{For $\beta^+$ decay $Q=(m_x-m_y-2m_e)c^2$}
\opt{D}{For $\beta^-$ decay $Q=(m_x-m_y+2m_e)c^2$}
\end{qoptions}
\qn{2} The potential difference between the target and the cathode of an X-ray tube is $50$ kV. The current in the tube is $20$ mA. Only $1\%$ of the total energy is emitted as X-rays. (Take: $hc=1240$ eV nm)
\begin{qoptions}
\opt{A}{The maximum frequency of the emitted x-rays is approx $1.2\times10^{19}$ Hz.}
\opt{B}{The heat should be removed at $900$W to keep the target at constant temperature}
\opt{C}{A wavelength of $0.2$ nm can be emitted by this tube}
\opt{D}{Electrons are being emitted at the rate of $1.25\times10^{17}$/sec by the cathode}
\end{qoptions}
\qn{3} A glass plate of refractive index $\mu_3$ is coated with a thin layer of thickness $t$ and refractive index $\mu_2$ ($\mu_2>\mu_3$). Light of wavelength $\lambda$ travelling in air is incident normally on layer. It is partially reflected at the upper and the lower surfaces of layer. Choose correct option(s)
\begin{qoptions}
\opt{A}{The two reflected waves interfere constructively if $t=\dfrac{\left(n+\frac12\right)\lambda}{2\mu_2}$.}
\opt{B}{The two reflected waves interfere constructively if $t=\dfrac{\left(n+\frac12\right)\lambda}{2\mu_3}$.}
\opt{C}{The two reflected waves interfere destructively if $t=\dfrac{n\lambda}{2\mu_2}$}
\opt{D}{The two reflected waves interfere destructively if $t=\dfrac{n\lambda}{2\mu_3}$}
\end{qoptions}
\qn{4} A wire under tension between two fixed point A \& B, executes sustained transverse vibrations so that the midpoint O of AB is a node, then
\qfig{p2_phy_s1_q04_1}
\begin{qoptions}
\opt{A}{There are no two points on the string which have a phase difference of $\dfrac\pi8$.}
\opt{B}{Total mechanical energy of a small element of string at O, remain constant.}
\opt{C}{Total mechanical energy of the string AB remain constant.}
\opt{D}{There can be only two points between A \& O, which will have equal speeds at all times.}
\end{qoptions}
\qn{5} Two mechanical waves, $y_1=2\sin2\pi(50\,t-2x)$ and $y_2=4\sin2\pi(ax+100\,t)$ propagate in a medium with same speed.
\begin{qoptions}
\opt{A}{The ratio of their intensities is $1:16$}
\opt{B}{The ratio of their intensities is $1:4$}
\opt{C}{The value of `$a$' is $4$ units}
\opt{D}{The value of `$a$' is $2$ units}
\end{qoptions}
\qn{6} A train, standing in a yard, blows a whistle of frequency $400$ Hz in still air. The wind starts blowing from yard to platform with a speed of $10$m/s. Given that the speed of sound in still air is $340$m/s,
\begin{qoptions}
\opt{A}{the frequency of sound as heard by an observer standing on the platform is $400$Hz.}
\opt{B}{the speed of sound for the observer standing on the platform is $350$m/s.}
\opt{C}{the frequency of sound as heard by the observer standing on the platform will increase.}
\opt{D}{the frequency of sound as heard by the observer standing on the platform will decrease.}
\end{qoptions}
\sect{SECTION-II}
\qn{1} A police car moving at $40$ m/s chases a thief running away at speed of $30$ m/s. The track is perpendicular to a stiff cliff as shown. The police man blows a horn at $40$ Hz. If sound has a speed of $340$ m/s, what is the beat frequency (in Hz) heard by the thief ?
\qfig{p2_phy_s2_q01_1}
\qn{2} A stationary tuning fork is in resonance with an air column in a pipe. If the tuning fork is moved with a speed of $2\ \mathrm{ms^{-1}}$ in front of the open end of the pipe and parallel to it, the length of the pipe should be changed for the resonance to occur with the moving tuning fork. If the speed of sound in air is $320\ \mathrm{ms^{-1}}$, the smallest value of the percentage change required in the length of the pipe is \underline{\hspace{1cm}}.
\qn{3} In the front of the upper slit of YDSE apparatus, a thin film of a liquid of refractive index $1.40$ is placed. It is a hot day the liquid starts evaporating from the surface. A beam of light of wavelength $560$ nm is incident onto the YDSE apparatus and the intensity $I$ at the centre of the screen is monitored. Figure gives intensity $I$ as a function of time $t$. The intensity changes because of evaporation from the two sides of the film. Assume that the film is flat and has parallel sides. Also assume that the film's thickness decreases at a constant rate. If the maximum intensity is $I_0$, then find $x$ if intensity at $t=10$ sec is $xI_0$.
\qfig{p2_phy_s2_q03_1}
\qn{4} The figure shows a wave propagating in a string. The speed of point A, which is $3$cm above the mean position is $7\sqrt3\pi$ cm/s upward at the moment shown. Find the speed (in m/s) of wave if the tension in the string is $4.9$N.
\qfig{p2_phy_s2_q04_1}
\qn{5} An interference is observed due to two coherent sources A and B having phase difference zero, separated by a distance $4\lambda$ along the $y$-axis where $\lambda$ is the wavelength of the source. If the distance of $3^{\mathrm{rd}}$ maxima on $x$-axis ($\Delta x=3\lambda$) from B is $n\lambda$. find the value of $n$.
\qfig{p2_phy_s2_q05_1}
\qn{6} A positron has $0.04$ MeV kinetic energy. When it collide with an free electron at rest then they annihilate each other. One of the two photon produced is found to be moving in direction of incident positron. It's energy $x$ MeV. Find $x$. Neglect any relativistic effect. [mass of electron = mass of positron $=\left(\dfrac{0.5}{C^2}\right)$ MeV]
\sect{SECTION-III}
\qn{1} In a fission reaction ${}^{236}_{92}\mathrm{U}\to{}^{117}X+{}^{117}Y+n+n$, the binding energy per nucleon of X and Y is $8.5$ MeV whereas of ${}^{236}\mathrm{U}$ is $7.58$ MeV. The total energy liberated is (approx.) $25N$ MeV. Fill the value of $N$
\qn{2} When thermal neutrons (negligible kinetic energy) are used to induce the reaction;

${}^{10}_{5}\mathrm{B}+{}^{1}_{0}n\longrightarrow{}^{7}_{3}\mathrm{Li}+{}^{4}_{2}\mathrm{He}$; $\alpha$ - particles are emitted with an energy of $1.83$ MeV.

Given the masses of boron, neutron \& $\mathrm{He}^4$ as $10.01167$, $1.00894$ \& $4.00386$ amu respectively. What is the mass of ${}^{7}_{3}\mathrm{Li}$ Assume that particles are free to move after the collision. Approximate your answer to nearest integer.
\qn{3} Flat aluminum electrode is illuminated by light with a wavelength of $83$ nm. What is the maximum distance (in mm) from the electrode surface will a photoelectron move, if outside the electrode there is a retarding electric field of $E=750$ V/m. Photoelectric threshold for aluminum is $332$ nm. If your answer is $X$ then fill the value of $X/3$. (Take $hc=1245$ eV nm)
\qn{4} A target element A is bombarded with electrons and the wavelengths of the characteristic spectrum are measured. A second characteristic spectrum is also obtained, because of an impurity in the target. The wavelengths of the $K_\alpha$ lines are $196$ pm (element A) and $169$ pm (impurity). If the atomic number of impurity is $z=(10\lambda-1)$. Find the value of $\lambda$. (atomic number of element A is 27).
\qfullbegin
\qn{5} The main scale of a vernier callipers reads in millimeter and its vernier is divided into $10$ divisions which coincides with $9$ divisions of the main scale. The reading for shown situation is found to be $\dfrac{23x}{10}$ mm. Find the value of $x$.
\qfig{p2_phy_s3_q05_1}
\qfullend
\qn{6} Intensity observed in an interference pattern is $I=I_0\sin^2\theta$. At $\theta=30^\circ$, intensity $I=5\pm0.0020$ W/m$^2$. If $I_0=20$ W/m$^2$ and percentage error in angle $\theta$ is $\dfrac{n}{\pi}\sqrt3\times10^{-2}\%$, then the value of $n$ is.
\partsec{2}{CHEMISTRY}
\sect{SECTION-I}
\qn{1} Select the correct statement(s) when $100$ ml of $0.1$ M $\mathrm{CH_3COOH}$ aqueous solution is mixed with following solutions at $25^\circ$C. [Assume volume of solid to be negligible]

[Given : $K_a\ \mathrm{CH_3COOH}=1\times10^{-5}$ M, $\log2=0.3$]
\begin{qoptions}
\opt{A}{On mixing with $50$ ml, $0.2$ M NaOH(aq), it will form a buffer solution having pH $=5$.}
\opt{B}{On mixing with $0.01$ mole of solid KOH, it will form an aqueous solution having pH $=9$}
\opt{C}{On mixing with $100$ ml, $0.1$ M HCl(aq), it will form an aqueous solution having pH $=1.3$}
\opt{D}{On mixing with $10$ ml, $0.2$ M NaOH(aq), it will form a buffer solution having pH $=4.6$.}
\end{qoptions}
\qn{2} A metal `M' crystallizes in two different crystal systems FCC and HCP. If $d$ = density, $z$ = effective number of atoms per unit cell and $[a,b,c]$ are edge lengths of the unit cell along the crystallographic axes $[x,y,z]$ then, select the correct option(s) :
\begin{qoptions}
\opt{A}{$\dfrac{d_{\mathrm{FCC}}}{d_{\mathrm{HCP}}}\neq1$}
\opt{B}{$\dfrac{a_{\mathrm{FCC}}}{c_{\mathrm{HCP}}}=\dfrac{\sqrt3}{2}$}
\opt{C}{$\dfrac{d_{\mathrm{FCC}}}{d_{\mathrm{HCP}}}=1$}
\opt{D}{$\dfrac{z_{\mathrm{FCC}}}{z_{\mathrm{HCP}}}=\dfrac23$}
\end{qoptions}
\qn{3} 
\qfig{p2_chem_s1_q03_1}
Choose the correct statement regarding above reaction sequence
\begin{qoptions}
\opt{A}{Compound (A) is stable due to resonance}
\opt{B}{Compound (C) gives positive test of unsaturation}
\opt{C}{Compound (B) gives effervescence of $\mathrm{CO_2}$ with $\mathrm{NaHCO_3}$}
\opt{D}{Compound (A) forms red colour dye with $\beta$-naphthol}
\end{qoptions}
\qn{4} Which of the following statements are correct?
\begin{qoptions}
\opt{A}{Decron is co polymer of ethyleneglycol terephthatic acid}
\opt{B}{Nylon-6 is homo polymer of caprolactum}
\opt{C}{Novalac is co polymer of phenol and formaldehyde}
\opt{D}{Buna-S is the co polymer of butadiene and styrene.}
\end{qoptions}
\qn{5} Fusion of $\mathrm{MnO_2}$ with KOH in presence of $\mathrm{O_2}$ produces a salt W. Alkaline solution of W upon electrolytic oxidation yields another salt X. The manganese containing ions present in W and X, respectively, are Y and Z. Correct statement(s) is (are)
\begin{qoptions}
\opt{A}{Y is diamagnetic in nature while Z is paramagnetic}
\opt{B}{Both Y and Z are coloured and have tetrahedral shape}
\opt{C}{In both Y and Z, hybridization is $d^3s$}
\opt{D}{In aqueous acidic solution, Y undergoes disproportionation reaction to give Z and $\mathrm{MnO_2}$.}
\end{qoptions}
\qn{6} Alkali metals are characterised by their :
\begin{qoptions}
\opt{A}{high electropositive character}
\opt{B}{high reduction potentials}
\opt{C}{low melting points}
\opt{D}{high solubility in liquid ammonia at $-33^\circ$C.}
\end{qoptions}
\sect{SECTION-II}
\qn{1} The enthalpy change for vapourisation of liquid `A' at $200$ K and $1$ atm is $25$ kcal/mol.

Given : $C_{p,m}$ (A, liquid) $=50$ Cal/mole,

$C_{p,m}$ (A, gas) $=40$ Cal/mole

Normal boiling point of liquid `A' is $300$ K.

[Given : $\ln2=0.7$ and $\ln3=1.1$]

Calculate change in Gibbs energy (in kcal/mole) for process A($\ell$, $200$ K $1$ atm) $\to$ A(g, $200$ K $1$ atm). Assume heat capacity to remain constant with temperature.
\qn{2} The enthalpy change for vapourisation of liquid `A' at $200$ K and $1$ atm is $25$ kcal/mol.

Given : $C_{p,m}$ (A, liquid) $=50$ Cal/mole,

$C_{p,m}$ (A, gas) $=40$ Cal/mole

Normal boiling point of liquid `A' is $300$ K.

$\ln2=0.7$ and $\ln3=1.1$

Calculate $\Delta S$ (in Cal/mole-K) for process A($\ell$, $200$ K $1$ atm) $\to$ A(g, $200$ K $1$ atm).
\qn{3} A tripeptide is composed equally of L-tyrosine, L-glycine and L-valine (one molecule of each). How many number of tripeptides can be obtained?
\qn{4} How many essential amino acids are present in the living system?
\qn{5} Total how many moles of NaOH are required to completely neutralise the solution formed by complete hydrolysis of $1$ mole of $\mathrm{BBr_3}$ ?
\qn{6} Consider the electronic configuration of lanthanides series. If the unpaired electrons in the ground state of $\mathrm{Ce^{+3}}$, $\mathrm{Yb^{+3}}$ and $\mathrm{Gd^{+3}}$ are $a$, $b$ and $c$ respectively, then the value of $a+b+c$ ?
\sect{SECTION-III}
\qn{1} An aqueous solution of KCl has a freezing point of $-0.51^\circ$C. Find the osmotic pressure (in atm) of same solution at $37^\circ$C. [Given : $K_f$ of $\mathrm{H_2O}(\ell)=1.85$ K-kg/mol, $R=0.082$ atm-L/mol-K]. Assume solution to be highly dilute.
\qn{2} Isohydric solutions are the solutions having the same concentration of hydronium ion. If $0.2$ M HA (monobasic acid) solution is isohydric with $4\times10^{-4}$ M HCl solution, then $K_a$ for HA is $P\times10^{-7}$. Value of `P' is :
\qn{3} Identify total number of Essential amino acids

(a) Glycine \quad (b) Alanine \quad (c) Phenyl alanine \quad (d) Arginine

(e) Cysteine \quad (f) Histidine \quad (g) Proline \quad (h) Tryptophan
\qn{4} How many of the following are condensation polymer?

(i) Nylon-6,6 \quad (ii) Buna-N \quad (iii) Nylon-6 \quad (iv) Terylene

(v) Teflon \quad (vi) Polystyrene \quad (vii) Glyptal \quad (viii) Bakelite
\qn{5} How many aqueous solution of compounds cause hardness in water?

$\mathrm{NaCl}$, $\mathrm{CaCl_2}$, $\mathrm{CaSO_4}$, $\mathrm{Mg(HCO_3)_2}$, $\mathrm{KNO_3}$, $\mathrm{Mg(NO_3)_2}$
\qn{6} The number of neutral colourless gases among the following is \underline{\hspace{0.8cm}}

$\mathrm{N_2O}$, NO, $\mathrm{N_2O_3}$, $\mathrm{NO_2}$, $\mathrm{N_2O_4}$, $\mathrm{N_2O_5}$
\partsec{3}{MATHEMATICS}
\sect{SECTION-I}
\qn{1} Let $A(z_1)$ be a point lies on the circle whose centre is at origin and radius is $4$ unit, $z_2$ be the midpoint of $O$ and $z_1$, where $O$ is origin. Let $\omega=e^{\frac{i2\pi}{21}}$ and $\displaystyle\sum_{r=1}^{20}|z_1+\omega^rz_2|^2=K$ then the sum of digit of $K$ is divisible by...
\begin{qoptions}
\opt{A}{$3$}
\opt{B}{$4$}
\opt{C}{$5$}
\opt{D}{$7$}
\end{qoptions}
\qn{2} Let A,B,C be $3$ sets such that

$A=\left\{(x,y):\dfrac{x}{\cos\theta}=\dfrac{y}{\sin\theta}=2,\theta\in R\right\}$

$B=\left\{(x,y):\dfrac{x+1}{\cos2205^\circ}=\dfrac{y}{\sin2205^\circ}=r,r\in R\right\}$

$C=\left\{(x,y):(x-\sqrt3)^2+(y-1)^2\leqslant a^2;a\in R\right\}$

Let $C_1,C_2$ be the curves represented by A and B respectively while R be the region represented by C then which of following statement is/are correct
\begin{qoptions}
\opt{A}{A points $P(x_1,y_1)$ and $Q(x_2,y_2)$ be such that $P,Q\in A\cap B$, then point of intersection of tangents drawn to $C_1$ at $P$ and $Q$ is $(2,3)$.}
\opt{B}{A points $P(x_1,y_1)$ and $Q(x_2,y_2)$ be such that $P,Q\in A\cap B$, then point of intersection of tangents drawn to $C_1$ at $P$ and $Q$ is $(-4,4)$.}
\opt{C}{Minimum value of `$a$' for which $A\cap C=A$ is $\alpha$, then length of curve $C_2$ lying in R corresponding to $a=\alpha$, is $\sqrt{58}$}
\opt{D}{Minimum value of `$a$' for which $A\cap C=A$ is $\alpha$, then length of curve $C_2$ lying in R corresponding to $a=\alpha$, is $\sqrt{57}$}
\end{qoptions}
\qn{3} If $x,y,z$ are non zero real numbers such that

$a^2+b^2+c^2+\dfrac{1}{x^2y^2z^2}\Big((x^2y+xy^2)^2$ $+(y^2z+zy^2)^2+(xz^2+z^2x)^2\Big)$
$=\dfrac{-2}{xyz}\left(ayz(y+z)+bxz(x+z)+cxy(x+y)\right)$

where $a,b,c\in R-\{1\}$ and the value of $\left[\left|\dfrac{1+a}{1-a}+\dfrac{1+b}{1-b}+\dfrac{1+c}{1-c}\right|\right]$ ($[.]$ greatest integer function) divides
\begin{qoptions}
\opt{A}{$2$}
\opt{B}{$3$}
\opt{C}{$7$}
\opt{D}{$5$}
\end{qoptions}
\qn{4} If
\par\noindent\hfill{\footnotesize$A=\begin{bmatrix}x+xy+xyz&0&0\\0&y+yz+xyz&0\\0&0&z+zx+xyz\end{bmatrix}$}\hfill\null\par
$|A|\neq0$ and $xyz=1$ and sum of maximum and minimum value of $\mathrm{Tr}(A^{-1})$ (where $\mathrm{Tr}(A)$ denotes sum of diagonal elements) is $\dfrac pq$, then which of following options is/are correct ($p$ and $q$ are coprime)
\begin{qoptions}
\opt{A}{$p+q=3$}
\opt{B}{$p+q=5$}
\opt{C}{$p-q=1$}
\opt{D}{$3p-2q=4$}
\end{qoptions}
\qn{5} Let $z$ be complex number such that $|z|=1$ then possible values of $|1-z|+|1+z^2|$ among the following is/are
\begin{qoptions}
\opt{A}{$\dfrac{1}{\sqrt2}$}
\opt{B}{$\sqrt2$}
\opt{C}{$\pi$}
\opt{D}{$4$}
\end{qoptions}
\qn{6} A box contains $n$ coins of which at least one is biased. Let $E_k$ denote the event that exactly $k$ out of the $n$ coins are biased. Let $P(E_k)$ is proportional to $k$ for $1\le k\le n$. If $A$ denotes the event that a coin selected out of $n$ coins at random, is biased, then :-
\begin{qoptions}
\opt{A}{$\displaystyle\lim_{n\to\infty}P(A)=\dfrac13$}
\opt{B}{$\displaystyle\lim_{n\to\infty}P(A)=\dfrac23$}
\opt{C}{$P\!\left(\dfrac{E_1}{A}\right)=\dfrac{3}{n(n+1)(2n+1)}$}
\opt{D}{$P\!\left(\dfrac{E_1}{A}\right)=\dfrac{6}{n(n+1)(2n+1)}$}
\end{qoptions}
\sect{SECTION-II}
\qn{1} If $a,b,c$ be three distinct natural numbers. Then the minimum value of
$\left|a+be^{i\frac{2\pi}{3}}+ce^{i\frac{4\pi}{3}}\right|^2+\left|a+be^{-i\frac{2\pi}{3}}+ce^{-i\frac{4\pi}{3}}\right|^2$ is
\qn{2} If $a,b,c$ be the distinct real numbers and $x,y,z$ be any three variables such that the system of linear equations

$(3z-2y)a+(2x-z)b+(y-3x)c=0$

$(y-3x)a+(3z-2y)b+(2x-z)c=0$

$(2x-z)a+(y-3x)b+(3z-2y)c=0$

can be express as $AX=O$, $A$ is $3\times3$ square matrix and $X=\begin{bmatrix}x\\y\\z\end{bmatrix}$ then the value of $\det(A)+\det(2A)+\det(3A)$ is :
\qn{3} The value of $\displaystyle\sum_{n=2}^{\infty}\dfrac{n^4+3n^2+10n+10}{2^n(n^4+4)}$ is
\qn{4} If $\alpha_1,\alpha_2,\alpha_3,\ldots,\alpha_{11}$ are roots of the equation $x^{11}+x^{10}+\ldots\ldots+x+1=0$ and the area of the polygon whose vertices are $1,\alpha_1,\alpha_2,\alpha_3,\ldots,\alpha_{11}$ is $\lambda$ square unit. Then the value of $\lambda$ is
\qn{5} Let $P$ be any point on ellipse $3x^2+4y^2=12$ and $S,S'$ are its foci then the locus of the centroid of triangle $PSS'$ is a conic $C$. Area of quadrilateral formed by tangents at the extremities of latus-rectum of the conic $C$ equals
\qn{6} The area of the region which consists of all the points satisfying the conditions $|x-y|+|x+y|\le8$, $xy\ge2$, $x\ge0$ and $y\ge0$ is equal to - (take $\ln2=0.69$)
\sect{SECTION-III}
\qn{1} Let $\vec a$ and $\vec b$ are two non-zero, non collinear vectors ($|\vec a|=1$) such that vectors $3(\vec a\times\vec b)$ and $2(\vec b-(\vec a.\vec b)\vec a)$ represents two sides of a triangle. If area of triangle is $\dfrac34\left(|\vec b|^4+4\right)$, then find the value of $|\vec b|^2$
\qn{2} For $f(x)=x^4+|x|$, Let $I_1=\displaystyle\int_0^{\pi}f(\cos x)dx$, $I_2=\displaystyle\int_0^{\pi/2}f(\sin x)dx$, then the value of $\dfrac{I_1}{I_2}$ is
\qn{3} If a normal chord AB to the hyperbola $xy=c^2$, normal at A, subtends an angle $\theta$ at centre and $\angle OAB=\alpha$, ($O$ is origin) then $\left|\dfrac{\tan\theta}{\tan\alpha}\right|$ is
\qn{4} Three numbers are chosen at random one after another with replacement, from the set $S=\{1,2,3,\ldots,100\}$. Let $P_1$ be the probability that the minimum of chosen numbers is at most $60$ and $P_2$ be the probability that the maximum of chosen numbers is at least $41$ and at most $80$. The value of $([P_1+P_2])$ is \underline{\hspace{0.8cm}}

(Where $[.]$ denotes greatest integer function)
\qn{5} A number is chosen at random from the set $\{1,2,3,\ldots,1000\}$. Let $P$ be the probability that the chosen number is a multiple of $3$ or a multiple of $7$. Then the value of $[5P]$ is \underline{\hspace{0.8cm}}

(Where $[.]$ denotes greatest integer function)
\qn{6} Let $M=[a_{ij}]_{2\times2}$ whose all entries are distinct and $a_{ij}\in\{1,2,5,10\}$, then number of such singular matrices $M$ is
\end{multicols}
\end{document}
